% ═══════════════════════════════════════════════════════════════════
% Section I: Introduction
% IEEE Communications Letters style (~450 words)
% v4: continuous alpha scheduling (two-mode)
% ═══════════════════════════════════════════════════════════════════

\section{Introduction}\label{sec:intro}

% ── 문단 1: CE default → cost 폭발 ───────────────────────────────
% NOTE (2026-03-25): "per-slot CE" 전제는 NR 표준 맥락에서만 정확.
% PACENet은 per-pilot-block, Channelformer는 per-frame CE. 학술 DL-CE 일반에는 해당 안 됨.
% 더 정확한 framing: "CE at every pilot observation interval" (per-slot이든 per-frame이든).
% 현재 문장은 NR DMRS 기반이므로 per-slot이 맞지만, 리뷰어가 DL-CE 논문을 근거로
% "per-slot이 default 아니다"라고 공격 가능. 필요시 아래 대안으로 교체:
%   "CE is conventionally executed at every pilot observation interval,
%    whether per-slot in NR or per-frame in other systems."
% 참고: docs/ce-skip/ce_execution_granularity.md
\IEEEPARstart{P}{er-slot} channel estimation (CE) has been an unquestioned default in OFDM systems.
3GPP NR maps the pilot signals, i.e., demodulation reference signals (DMRS), to every scheduled slot, and the channel learned from DMRS is assumed valid only within that slot~\cite{3gpp_ts38211}.
Receivers therefore conventionally refresh CE on a per-slot basis~\cite{edfors1998ofdm}.
When CE overhead was negligible, this convention posed no problem.
However, sixth-generation extremely large aperture arrays (6G ELAA) scale antennas to hundreds or thousands of elements~\cite{cui2022channel, bjornson2024towards}, and the cost of advanced CE methods grows rapidly with array size. Linear minimum mean-square error (LMMSE) estimation requires matrix inversion that scales cubically in antenna dimension~\cite{shariati2014low}, while DL-based CE (DL-CE) adds neural-network inference overhead~\cite{requestnet2025, channelformer2023}.

% ── 문단 2: prior work (per-invocation만) ────────────────────────
Existing approaches tackle this growing overhead in various ways.
Early-exit networks reduce per-inference cost by terminating iterative refinement early~\cite{icenet2025}, model compression shrinks estimator complexity~\cite{channelformer2023}, and temporal prediction replaces CE with learned channel forecasting~\cite{continual_pred2025}.
% NOTE: "every slot"은 NR 기준. DL-CE 논문들은 per-frame/per-block이므로 "every observation"이 더 정확.
Yet all of these still require a forward pass at every observation interval.
Pilot overhead reduction minimizes cost within a single CE invocation~\cite{channelformer2023}, and pilot skipping exploits coherence time to omit pilot transmission for slow users~\cite{abboud2017channel}. Neither addresses whether to \emph{execute} the CE algorithm on already-received pilots.


% ── 문단 3: enabler (기회 포착) ──────────────────────────────────
We observe that a complementary opportunity has emerged from the architectural side.
In traditional fixed-function base stations, baseband processing runs on dedicated ASIC/DSP hardware~\cite{park2025vran}, where CE executes as part of a rigid pipeline that cannot be conditionally bypassed~\cite{venkataramani2020spectrum}.
GPU-native platforms have removed this constraint by deploying CE as a CUDA kernel~\cite{nvidia_aerial}, and O-RAN dApps enable per-slot PHY control at sub-millisecond latency~\cite{lacava2025dapps, polese2025airan}.
Rather than reducing per-invocation cost, it is now possible to reduce invocation frequency itself systematically.

% ── 문단 4: key insight → tradeoff → contributions ───────────────
Based on this observation, we shift the focus from \emph{how} to estimate to \emph{when} to estimate, and propose \emph{opportunistic CE inference scheduling}.
The key insight is that low-cost mandatory pilot reception, such as least-squares (LS) estimation, and orders-of-magnitude costlier discretionary full CE inference, such as LMMSE or DL-CE, are separable.
The former can monitor channel variation at each slot, while the latter need only be launched when reusing the previous estimate would cause non-negligible degradation.
Skipping CE saves computation but leaves the receiver with an aged channel estimate, degrading beamforming gain and equalization quality~\cite{truong2013effects}.
A practical scheduling policy must balance these savings against quality loss, keeping NMSE degradation within an acceptable bound.

Our Sionna~RT-based evaluations on 5G massive MIMO ($8\!\times\!8$, 3.5\,GHz) and 6G ELAA ($24\!\times\!24$, 15/28\,GHz) channels in a Munich urban scenario results in ...XXX.


To the best of our knowledge, this is the first work to formalize CE inference scheduling as a decision layer orthogonal to CE algorithm design.
We make two contributions.

\begin{enumerate}
\item We formalize \emph{CE inference scheduling} as a per-slot decision problem that decouples mandatory pilot reception from discretionary CE inference, and propose a two-mode adaptive scheduling scheme. In \emph{Blend} mode, the estimate is updated via exponential smoothing whose weight varies continuously with observed channel variation. In \emph{Full} mode, triggered when variation exceeds a threshold, the complete CE algorithm is launched.

\item We demonstrate that the proposed scheme is CE-algorithm-agnostic: it yields consistent computation\,--\,accuracy tradeoffs across LS, genie-aided LMMSE, and DL-CE without method-specific tuning, and scheduling thresholds calibrated on a subset of base stations transfer to unseen sites with ${<}\,2$\,dB NMSE degradation.
\end{enumerate}

The paper is organized as follows.